% FILE: 04_implementation_surface.tex
% Project: adversarial
% Thesis Role: adversarial-testing-and-hostile-assumption-verification

\section{Implementation Surface}
\label{sec:adversarial-implementation}

\subsection{Source Files}
\label{sec:adversarial-implementation-sources}

The adversarial implementation surface consists of two executable Python scripts and twelve
supporting research documents.

\textbf{\texttt{adversarial/simulate\_laundering\_spoofing.py}} (531 lines) is the primary
simulation engine. It implements the \texttt{IngestionBoundary}, \texttt{LaunderingRefusalError},
and \texttt{Wasm4pmBridge} classes and the \texttt{run\_simulation\_tests} function containing
ten scenarios with Python \texttt{assert} statements. The script requires Python 3 with
\texttt{json}, \texttt{hashlib}, \texttt{hmac}, \texttt{uuid}, \texttt{ctypes}, and
\texttt{pandas}. Scenarios 0 through 5 operate entirely in Python and do not require
\texttt{libwasm4pm.dylib}. Scenarios 6 through 9 require the compiled library at
\texttt{sources/wasm4pm/target/debug/libwasm4pm.dylib} or the equivalent \texttt{.so} path;
the \texttt{locate\_library()} function probes four candidate paths and raises
\texttt{FileNotFoundError} if none exists.

\textbf{\texttt{adversarial/validate\_standards.py}} (139 lines) validates six standards
placement files: XES (\texttt{xes\_process-intelligence\_placement.md}), OCEL 2.0
(\texttt{ocel\_process-intelligence\_placement.md}), BPMN
(\texttt{bpmn\_process-intelligence\_placement.md}), POWL
(\texttt{powl\_placement.md}), Declare (\texttt{declare\_placement.md}), and OCPQ
(\texttt{ocpq\_placement.md}). For each placement file the script applies four checks:
placeholder term detection (regex against TODO, FIXME, unimplemented, stub, tbd),
broken-link validation for \texttt{file:///} URIs against the filesystem, malformed JSON
block detection, and presence of the Trans-Standard Conversions section with
\texttt{LossReport} schema and \texttt{witness\_signature} field. The script exits 0 on
pass and 1 on fail; both exit codes are meaningful for CI pipeline integration.

Supporting research documents include \texttt{adversarial/RED\_TEAM\_SCOPE.md} (charter),
\texttt{adversarial/RED\_TEAM\_FINDINGS\_001.md} and
\texttt{adversarial/RED\_TEAM\_FINDINGS\_002.md} (ten structured findings),
\texttt{adversarial/simulation\_models\_specification.md} (verified execution logs for six
primary scenarios), \texttt{standards/ocel-adversarial-gravity.md} (PROV-O gap analysis),
\texttt{sources/papers/adversarial-canon.md} (Sybil-attack canon),
\texttt{sources/wasm4pm-compat/adversarial-type-law.md} (three structural laws for witness
lattice integrity), \texttt{experiments/raw-laundering\_refusal\_sample.md} (concrete JSON
refusal instance), \texttt{prompts/execution-plans/raw-laundering-block.md} (downstream
execution plan), and \texttt{atlas/agi\_adversarial\_defenses.md} (AGI adversarial defense
atlas).

\subsection{Commands}
\label{sec:adversarial-implementation-commands}

The simulation engine is invoked directly:
\begin{verbatim}
python3 adversarial/simulate_laundering_spoofing.py
\end{verbatim}
A successful run prints scenario-by-scenario verdicts and terminates with the message
\texttt{ALL ADVERSARIAL SIMULATION SCENARIOS SUCCESSFULLY VERIFIED}. Scenarios 6 through 9
will raise \texttt{FileNotFoundError} if \texttt{libwasm4pm.dylib} is absent; this is a
known open question (see Section~\ref{sec:adversarial-questions-gaps}).

The standards audit script is invoked from the repository root:
\begin{verbatim}
python3 adversarial/validate_standards.py
\end{verbatim}
A passing audit prints \texttt{AUDIT PASSED} and exits 0. A failing audit prints the
per-file error list and exits 1.

\subsection{Data and Ontology Surfaces}
\label{sec:adversarial-implementation-data}

The adversarial corpus interacts with three data surfaces. The \textbf{log envelope format}
is a Python dictionary with a \texttt{traces} list, where each trace contains a
\texttt{trace\_id} and an \texttt{events} list. Each event contains a \texttt{payload}
dictionary (with \texttt{event\_id}, \texttt{activity\_name}, \texttt{timestamp\_ns},
\texttt{attributes}), a \texttt{signature} (HMAC-SHA256 hexdigest), and a
\texttt{hash\_chain\_link} (SHA-256 hexdigest). Canonical serialization uses
\texttt{json.dumps} with \texttt{sort\_keys=True} and \texttt{separators=(',', ':')}.

The \textbf{wasm4pm linear memory layout} is managed via \texttt{wasm\_alloc} offsets.
Heap reads and writes use \texttt{ctypes.memmove} through absolute pointers returned by
\texttt{wasm\_get\_absolute\_ptr}. The \texttt{wasm\_shred\_heap} function applies a
four-pass zeroization protocol; Scenario 9 verifies that sentinel bytes are unreadable after
the call.

The \textbf{OCEL adversarial ontology surface} is analyzed in
\texttt{standards/ocel-adversarial-gravity.md}: the OCEL SQLite schema (\texttt{event},
\texttt{object}, \texttt{event\_object} tables) and the OCEL JSON schema (\texttt{events},
\texttt{objects}, \texttt{relationships} arrays) both lack PROV-O enforcement. The Reality
Receipt mechanism addresses this by wrapping OCEL/XES events in a signed cryptographic
envelope at the wasm4pm interception proxy.

\subsection{Templates and Rendered Artifacts}
\label{sec:adversarial-implementation-templates}

The primary rendered artifact is the \texttt{LogLaunderingRefusalLog} JSON schema, whose
concrete instance is documented in \texttt{experiments/raw-laundering\_refusal\_sample.md}.
This schema specifies the structure of refusal responses: \texttt{evaluation\_id},
\texttt{analyzed\_log\_hash\_sha256}, \texttt{verdict}, \texttt{detection\_rules} (a boolean
map of all five rules), and \texttt{refusal\_reason} (\texttt{rule\_failed},
\texttt{trace\_id}, \texttt{detailed\_error}). The schema serves as the contract between the
ingestion boundary and any downstream audit or gap-remediation workflow.

The \texttt{RED\_TEAM\_FINDINGS} template requires each finding to specify: finding number,
severity (CRITICAL/MAJOR/MINOR/EXPECTED), challenge category from the seven-category charter,
specific document and claim challenged, and remediation path or the token
\texttt{[UNRESOLVABLE PRE-ALIVE]}.
